% \begin{document}
\section{Reference Types}

We extend the language with OCaml-style mutable references. The syntax and
semantics are:
\begin{align*}
e &::= \cdots \mid \ttkw{ref}\ e \mid\, !e \mid e_1 := e_2 \mid \ell \\
v &::= \cdots \mid \ell \\
E &::= \cdots \mid \ttkw{ref}\ E \mid\, !E \mid E := e \mid v := E
\end{align*}
\[
\inferrule*[Right=\textsc{Alloc}]
  {\ell \notin \dom(\sigma)}
  {\langle \ttkw{ref}\ v, \sigma\rangle \to \langle \ell, \sigma[\ell \mapsto v]\rangle}
\qquad
\inferrule*[Right=\textsc{Deref}]
  {\sigma(\ell) = v}
  {\langle !\ell, \sigma\rangle \to \langle v, \sigma\rangle}
\]
\[
\inferrule*[Right=\textsc{Assign}]
  {\ell \in \dom(\sigma)}
  {\langle \ell := v, \sigma\rangle \to \langle (), \sigma[\ell \mapsto v]\rangle}
\]

We add a type $\tau\ \kw{ref}$ for locations holding values of type $\tau$. For
instance, $\ttkw{ref}\ 7 : \kw{int\ ref}$. Assignment is performed only for its
effect, so it has type $\kw{unit}$.
\[
\tau ::= \cdots \mid \tau\ \kw{ref}
\]
\[
\inferrule{\Gamma \vdash e : \tau}{\Gamma \vdash \ttkw{ref}\ e : \tau\ \kw{ref}}
\qquad
\inferrule{\Gamma \vdash e : \tau\ \kw{ref}}{\Gamma \vdash !e : \tau}
\qquad
\inferrule{\Gamma \vdash e_1 : \tau\ \kw{ref} \\ \Gamma \vdash e_2 : \tau}
          {\Gamma \vdash e_1 := e_2 : \kw{unit}}
\]

\subsection*{Store typings}

There is no rule yet for typing a location $\ell$. Its type should be
$\tau\ \kw{ref}$ where $\tau$ is the type of the value stored at $\ell$, but
inspecting the store is inefficient and may be inconclusive (e.g.\ if
$\sigma(\ell) = \ell$). Instead we use \emph{store typings} $\Sigma$, which are
partial functions from locations to types. The typing judgment now has four
components and is written $\Gamma, \Sigma \vdash e : \tau$. The reference rules
become (other rules are modified to carry $\Sigma$ in the obvious way):
\[
\inferrule{\Gamma, \Sigma \vdash e : \tau}{\Gamma, \Sigma \vdash \ttkw{ref}\ e : \tau\ \kw{ref}}
\qquad
\inferrule{\Gamma, \Sigma \vdash e : \tau\ \kw{ref}}{\Gamma, \Sigma \vdash !e : \tau}
\]
\[
\inferrule{\Gamma, \Sigma \vdash e_1 : \tau\ \kw{ref} \\ \Gamma, \Sigma \vdash e_2 : \tau}
          {\Gamma, \Sigma \vdash e_1 := e_2 : \kw{unit}}
\qquad
\inferrule{\Sigma(\ell) = \tau}{\Gamma, \Sigma \vdash \ell : \tau\ \kw{ref}}
\]

\subsection*{Type soundness}

The earlier soundness statement must be adapted to account for the store and the
store typing.

\medskip
\noindent\textbf{Definition.} A store $\sigma$ is \emph{well-typed} with respect to
$\Gamma$ and $\Sigma$, written $\Gamma, \Sigma \vdash \sigma$, if
$\dom(\sigma) = \dom(\Sigma)$ and for all $\ell \in \dom(\sigma)$ we have
$\Gamma, \Sigma \vdash \sigma(\ell) : \Sigma(\ell)$.

\thm{Type soundness}{
If $\cdot, \Sigma \vdash e : \tau$ and $\cdot, \Sigma \vdash \sigma$ and
$\langle e, \sigma\rangle \to^* \langle e', \sigma'\rangle$, then either $e'$ is a
value, or there exist $e''$ and $\sigma''$ such that
$\langle e', \sigma'\rangle \to \langle e'', \sigma''\rangle$.
}

The proof uses progress and preservation as before. Progress adapts easily.
Preservation is subtler: \textsc{Alloc} grows the store, so $\dom(\sigma') \neq
\dom(\Sigma)$, and the store typing must be allowed to grow too.

\mlenma{Preservation}{
If $\Gamma, \Sigma \vdash e : \tau$ and $\Gamma, \Sigma \vdash \sigma$ and
$\langle e, \sigma\rangle \to \langle e', \sigma'\rangle$, then there exists
$\Sigma' \supseteq \Sigma$ such that $\Gamma, \Sigma' \vdash e' : \tau$ and
$\Gamma, \Sigma' \vdash \sigma'$.
}

Here $\Sigma' \supseteq \Sigma$ means $\Sigma(\ell) = \Sigma'(\ell)$ for all
$\ell \in \dom(\Sigma)$; viewing $\Sigma$ as the set of pairs
$\{(\ell, \tau) \mid \ell \in \dom(\Sigma) \wedge \Sigma(\ell) = \tau\}$, this is
set inclusion. Intuitively, $\Sigma'$ is either $\Sigma$ or $\Sigma$ extended with
a newly allocated location.

\subsection*{Landin's knot}

References recover recursion (and hence Turing completeness): create a cell
holding a dummy function and ``backpatch'' it with the real definition.
\begin{verbatim}
let r = ref (fun x:int. 0) in
let f = fun x:int. if x = 0 then 1 else x * (!r) (x - 1) in
let a = (r := f) in
f 5
\end{verbatim}
After the assignment, $!r$ is $f$ itself, so $f\ 5$ evaluates to $120$.

\section{Recursive Types}

Languages commonly allow recursive data types, e.g.\ in OCaml
\begin{verbatim}
type tree = Leaf | Node of int * tree * tree
\end{verbatim}
We want $\mathit{tree}$ to satisfy
$\mathit{tree} = \kw{unit} + \kw{int} \times \mathit{tree} \times \mathit{tree}$,
i.e.\ to be a solution of
\[
\alpha = \kw{unit} + \kw{int} \times \alpha \times \alpha,
\]
but no such type exists among the types seen so far. There are two standard
approaches: \emph{equirecursive} and \emph{isorecursive}.

\section{Equirecursive Types}

Repeatedly unwinding the equation gives
\begin{align*}
\alpha &= \kw{unit} + \kw{int} \times \alpha \times \alpha \\
       &= \kw{unit} + \kw{int} \times (\kw{unit} + \kw{int} \times \alpha \times \alpha)
          \times (\kw{unit} + \kw{int} \times \alpha \times \alpha) \\
       &= \cdots
\end{align*}
a sequence of ever deeper finite trees, each a substitution instance of the
previous one. Taking the limit yields an infinite labeled tree whose nodes are
labeled by the constructors $\times, +, \kw{int}, \kw{unit}$; this infinite tree
is the solution, and it is what we take as the type $\mathit{tree}$.

Finite types over constructors such as $\to, \times, +, \kw{unit}, \kw{int}$ are
finite labeled trees (e.g.\ $\kw{int} \to \kw{int} \to \kw{int}$). Infinite types
are infinite labeled trees respecting arities: binary constructors have exactly
two children, nullary ones are leaves.

\begin{center}
\begin{tikzpicture}[level distance=8mm, sibling distance=14mm]
\node {$\to$}
  child {node {$\kw{int}$}}
  child {node {$\to$}
    child {node {$\kw{int}$}}
    child {node {$\to$}
      child {node {$\kw{int}$}}
      child {node {$\vdots$}}}};
\end{tikzpicture}
\end{center}

A (finite or infinite) tree with only finitely many distinct subtrees (up to
isomorphism) is called \emph{regular}. The tree above is regular: its only
subtrees are itself and $\kw{int}$. The limit defining $\mathit{tree}$ is also
regular, with five subtrees: $\mathit{tree}$, $\kw{unit}$,
$\kw{int} \times \mathit{tree} \times \mathit{tree}$,
$\mathit{tree} \times \mathit{tree}$, and $\kw{int}$.

Regular trees suffice for finite systems of equations
\[
\alpha_1 = \tau_1, \quad \ldots, \quad \alpha_n = \tau_n,
\]
where each $\tau_i$ is a finite type over the constructors and
$\alpha_1, \ldots, \alpha_n$. Such a system has a solution $\sigma_1, \ldots,
\sigma_n$ of regular trees, and the solution is unique if no $\tau_i$ is a bare
variable.

\subsection{The $\mu$ constructor}

The solution of $\alpha = \tau$ (with $\tau \neq \alpha$) is finite if $\alpha
\notin \mathit{FV}(\tau)$ and infinite otherwise; we denote it $\mu\alpha.\,\tau$.
The operator $\mu$ binds $\alpha$ just as $\lambda$ binds variables in terms
(same scope, free/bound variables, $\alpha$-conversion, substitution). Since
$\mu\alpha.\,\tau$ solves $\alpha = \tau$,
\[
\mu\alpha.\,\tau = \tsub{\tau}{\mu\alpha.\,\tau}{\alpha}.
\]
For example,
\[
\mathit{tree} \triangleq \mu\alpha.\ \kw{unit} + \kw{int} \times \alpha \times \alpha.
\]
Solutions of mutually recursive systems are also expressible with $\mu$. If
neither $\tau_1$ nor $\tau_2$ is a variable, the system
$\alpha_1 = \tau_1,\ \alpha_2 = \tau_2$ has the unique solution
\[
\sigma_1 = \mu\alpha_1.\ \tsub{\tau_1}{\mu\alpha_2.\,\tau_2}{\alpha_2}
\qquad
\sigma_2 = \mu\alpha_2.\ \tsub{\tau_2}{\mu\alpha_1.\,\tau_1}{\alpha_1}.
\]
Mutual recursion arises in practice, e.g.\ Java classes \texttt{Node} and
\texttt{Edge} that refer to each other.

\subsection{Typing rules}

Since $\mu\alpha.\,\tau = \tsub{\tau}{\mu\alpha.\,\tau}{\alpha}$, the rules are
simple (equivalently: allow substituting equals for equals in types):
\[
\inferrule*[Right=$\mu$-\textsc{Intro}]
  {\Gamma \vdash e : \tsub{\tau}{\mu\alpha.\,\tau}{\alpha}}
  {\Gamma \vdash e : \mu\alpha.\,\tau}
\qquad
\inferrule*[Right=$\mu$-\textsc{Elim}]
  {\Gamma \vdash e : \mu\alpha.\,\tau}
  {\Gamma \vdash e : \tsub{\tau}{\mu\alpha.\,\tau}{\alpha}}
\]

\section{Isorecursive Types}

Here there are no infinite types: $\mu\alpha.\,\tau$ is itself a type, and it is
\emph{distinct but isomorphic} to $\tsub{\tau}{\mu\alpha.\,\tau}{\alpha}$.
Substituting $\mu\alpha.\,\tau$ for $\alpha$ is \emph{unfolding}; the reverse is
\emph{folding}, and both are explicit operations:
\begin{align*}
\ttkw{unfold}_{\mu\alpha.\,\tau} &: \mu\alpha.\,\tau \to \tsub{\tau}{\mu\alpha.\,\tau}{\alpha} \\
\ttkw{fold}_{\mu\alpha.\,\tau} &: \tsub{\tau}{\mu\alpha.\,\tau}{\alpha} \to \mu\alpha.\,\tau
\end{align*}
Subscripts are often suppressed.

\subsection{Typing rules}
\[
\inferrule*[Right=$\mu$-\textsc{Intro}]
  {\Gamma \vdash e : \tsub{\tau}{\mu\alpha.\,\tau}{\alpha}}
  {\Gamma \vdash \ttkw{fold}\ e : \mu\alpha.\,\tau}
\qquad
\inferrule*[Right=$\mu$-\textsc{Elim}]
  {\Gamma \vdash e : \mu\alpha.\,\tau}
  {\Gamma \vdash \ttkw{unfold}\ e : \tsub{\tau}{\mu\alpha.\,\tau}{\alpha}}
\]

\subsection{Operational semantics}

Under CBV, $\ttkw{fold}\ v$ is a value, and arguments of $\ttkw{fold}$ and
$\ttkw{unfold}$ are evaluated first:
\begin{align*}
v &::= \cdots \mid \ttkw{fold}\ v \\
E &::= \cdots \mid \ttkw{fold}\ E \mid \ttkw{unfold}\ E
\end{align*}
with the single new rule
\[
\ttkw{unfold}\ (\ttkw{fold}\ v) \to v.
\]
A value of type $\mu\alpha.\,\tau$ can only have been created by a $\ttkw{fold}$,
so unfolding it always succeeds.

\subsection{Example}

Let $\mathit{intlist} \triangleq \mu\alpha.\ \kw{unit} + \kw{int} \times \alpha$.
To sum a list $l : \mathit{intlist}$ with a recursive function $f$, we must
unfold $l$ to expose the sum type before casing:
\[
\ttkw{case}\ (\ttkw{unfold}\ l)\ \ttkw{of}\ 
(\lambda u : \kw{unit}.\ 0)\ \mid\ 
(\lambda p : \kw{int} \times \mathit{intlist}.\ (\#1\ p) + f\ (\#2\ p)).
\]
This is the OCaml code with the folds and unfolds that OCaml hides made explicit.

\section{Equirecursive vs.\ Isorecursive}

Java and Modula-3 are equirecursive: folded and unfolded types are equal and
$\ttkw{fold}$/$\ttkw{unfold}$ act as the identity, so e.g.\ a class
\texttt{E} with a field \texttt{E e} permits \texttt{e.e.e.e.x}. The ML family,
CLU, and C are isorecursive: $\mu\alpha.\,\tau$ and
$\tsub{\tau}{\mu\alpha.\,\tau}{\alpha}$ are different types and explicit casts
are needed (CLU calls them ``up'' and ``down''). In OCaml, \texttt{match} and
\texttt{let} perform $\ttkw{unfold}$ implicitly, and datatype constructors act as
$\ttkw{fold}$:
\begin{verbatim}
# type tree = Leaf | Node of int * tree * tree;;
# Node (4, Leaf, Leaf);;
- : tree = Node (4, Leaf, Leaf)
\end{verbatim}

\section{Numbers as a Recursive Type}

$\kw{bool}$ can be encoded as $\kw{unit} + \kw{unit}$ (with $\mathit{true}$ and
$\mathit{false}$ the left and right injections). With recursive types, $\kw{int}$
need not be primitive either. Natural numbers are
\begin{align*}
\mathit{nat} &\triangleq \mu\alpha.\ \kw{unit} + \alpha \\
0 &\triangleq \ttkw{fold}\ (\ttkw{inl}_{\kw{unit}+\mathit{nat}}\ ()) \\
1 &\triangleq \ttkw{fold}\ (\ttkw{inr}_{\kw{unit}+\mathit{nat}}\ 0) \\
2 &\triangleq \ttkw{fold}\ (\ttkw{inr}_{\kw{unit}+\mathit{nat}}\ 1),
\end{align*}
and so on. Arithmetic can be coded over $\mathit{nat}$, e.g.\ successor:
\[
(\lambda x : \mathit{nat}.\ \ttkw{fold}\ (\ttkw{inr}_{\kw{unit}+\mathit{nat}}\ x)) : \mathit{nat} \to \mathit{nat}.
\]
Thus $\kw{unit}$, functions, products, sums, and recursive types suffice to build
integers, lists, trees, floats, etc.

\section{Self-Application and $\Omega$}

Recall $\omega \triangleq \lambda x.\, x\, x$ and $\Omega \triangleq \omega\,\omega$.
In $\omega$, $x$ is applied as a function, so $x : \sigma \to \tau$, but it is also
its own argument, so $\sigma = \sigma \to \tau$. This suggests
$\mu\alpha.\,(\alpha \to \tau)$ for any $\tau$. Unfolding gives
\[
\ttkw{unfold}\ x : (\mu\alpha.\,(\alpha \to \tau)) \to \tau,
\]
so $(\ttkw{unfold}\ x)\ x : \tau$ and
\[
\omega \triangleq (\lambda x : \mu\alpha.\,(\alpha \to \tau).\ (\ttkw{unfold}\ x)\ x)
  : (\mu\alpha.\,(\alpha \to \tau)) \to \tau.
\]
Folding gives $\ttkw{fold}\ \omega : \mu\alpha.\,(\alpha \to \tau)$, hence
\[
\omega\ (\ttkw{fold}\ \omega) : \tau,
\]
which is $\Omega$ with explicit folds and unfolds. In OCaml:
\begin{verbatim}
# type u = Fold of (u -> u);;
# let omega = fun x -> match x with Fold f -> f x;;
val omega : u -> u = <fun>
# omega (Fold omega);;   (* runs forever *)
\end{verbatim}
So recursive types alone, without $\ttkw{fix}$ or references, allow
non-terminating yet well-typed programs.

\section{Untyped to Typed $\lambda$-Calculus}

With recursive types every pure untyped $\lambda$-term can be typed by
\[
U \triangleq \mu\alpha.\ \alpha \to \alpha,
\]
via the translation
\begin{align*}
[\![x]\!] &\triangleq x \\
[\![e_0\ e_1]\!] &\triangleq (\ttkw{unfold}\ [\![e_0]\!])\ [\![e_1]\!] \\
[\![\lambda x.\, e]\!] &\triangleq \ttkw{fold}\ (\lambda x : U.\ [\![e]\!]).
\end{align*}
Every untyped term maps to a term of type $U$.
% \end{document}
